%% P06 --- Matrices as linear maps
%%
%% Every number in this program that the reader cannot do in their head comes
%% from code/p06_matrices_as_maps.py and is pulled in with \val{}. The console
%% listing is a committed transcript written by that script.
%%
%% THE THESIS: a matrix is a function that respects addition and scaling, and
%% matrix multiplication is composition of those functions -- which is why it
%% is associative, why it does not commute, and why shape is a type signature.
%%
%% WHAT THIS PROGRAM IS OWED, read out of the files rather than remembered:
%%   F05  ALREADY DERIVES THE COLLAPSE, in one dimension and in full, including
%%        the relu argument -- two hundred linear layers collapse to one, and
%%        the activation exists because without it the composition is provably
%%        a waste. So the brief's promise to derive it again is SPENT.
%%   F06  names what is actually left, and it is the right job: "P06 makes the
%%        weights a matrix and shows that matrix multiplication is composition,
%%        which is F05's collapse argument done in more than one dimension."
%%   P03  needs shapes and what a matrix is before a transformer block can be
%%        counted; the count itself it hands to P32.
%%   P04  a rotation is itself a matrix, so "does this survive a change of
%%        basis" is a question about this program's object.
%%   P05  "P06 adds the next object, and the pattern repeats."
%%
%% WHAT THIS PROGRAM LEAVES ALONE, checked against tools/programs.json. The
%% brief promises multi-head attention as a reshape and a batch as one extra
%% index; BOTH ARE P07'S, which owns index notation, einsum, broadcasting,
%% reshape against transpose against permute, and the axes of a rank-4
%% attention tensor. Writing them here would spend P07's payoff.
%%   index notation, einsum, broadcasting, reshape/transpose/permute   -> P07
%%   rank, the four subspaces, least squares, LoRA                     -> P08
%%   the determinant, the inverse, change of basis                     -> P09
%%   counting a transformer block's parameters                         -> P32
%%
%% Twin: programs/pl/P06-matrices-as-maps.tex. Frame for frame, macro for
%% macro, number for number; tools/parity.py compares them.

\program{Matrices as linear maps}
\label{prog:P06}
\index{matrix!as a linear map}
\index{matrix!multiplication}

\begin{outcomes}
\outcome{1--8}{say what makes a function linear, and test a given rule against
  the definition}
\outcome{9--14}{multiply two matrices as a composition of functions rather than
  as a rule about rows and columns}
\outcome{15--19}{say why matrix multiplication does not commute, and give a
  pair of maps whose two orders disagree}
\outcome{20--25}{read a shape as a type signature, and say which compositions
  are legal before computing anything}
\outcome{26--36}{collapse a stack of linear layers to one, price what the stack
  bought, and say what the activation is for}
\end{outcomes}

\begin{quiz}
\item What two things must a function do to be called linear? \teachesat{2--3}
  \answerto{Respect addition and respect scaling: $f(u + v) = f(u) + f(v)$ and
  $f(cv) = c\,f(v)$. Everything else in this program follows from those two.}
\item Is $f(x) = x + 3$ linear? \teachesat{4--5}
  \answerto{No, despite being a straight line. $f(0) = 3$, and a linear
  function must send the zero vector to the zero vector. It is
  \emph{affine} \dash{} linear plus a shift.}
\item What does the product $AB$ mean, if $A$ and $B$ are functions?
  \teachesat{9--10}
  \answerto{Do $B$ first, then $A$. The product is the composition, and the
  right-to-left order is why it looks backwards.}
\item Give two operations on the plane whose two orders disagree.
  \teachesat{15--16}
  \answerto{A quarter turn and a flattening onto the horizontal axis. Project
  then rotate sends $(1, 0)$ to $(\val{p06.order.rp.x},
  \val{p06.order.rp.y})$; rotate then project sends it to
  $(\val{p06.order.pr.x}, \val{p06.order.pr.y})$.}
\item An $m \times n$ matrix multiplies an $n \times p$ matrix. What is the
  shape of the answer, and what does that shape say? \teachesat{20--21}
  \answerto{$m \times p$: a map from a space of $p$ dimensions to a space of
  $m$. The inner $n$ has to match and then disappears, which is exactly what a
  type signature does at a function call.}
\item Two linear layers with nothing between them. How many layers is that,
  really? \teachesat{26--27}
  \answerto{One. $W_2(W_1x + b_1) + b_2 = (W_2W_1)x + (W_2b_1 + b_2)$, which is
  a single matrix and a single bias. Program~\ref{prog:F05} showed this in one
  dimension; the matrices change nothing about the argument.}
\item $A$, $B$ and $C$ are matrices. Does it matter how you bracket $ABC$?
  \teachesat{33--34}
  \answerto{Not for the answer \dash{} it is associative, and that is a
  theorem. It matters enormously for the work: two bracketings of one product
  here cost $\val{p06.cost.left}$ and $\val{p06.cost.right}$ multiplications.}
\end{quiz}

% ============================================================
\section{What makes a function linear}
\label{sec:P06-linear}
\index{matrix!as a linear map}

\begin{fr}
Program~\ref{prog:P05} closed by saying that one operation added to a bare
space bought lengths, angles, perpendicularity and projection. This program
adds an object rather than an operation, and the same thing happens.

Here is the object as everybody first meets it: a rectangle of numbers.

\[ A = \begin{pmatrix} 2 & 0 \\ 0 & 3 \end{pmatrix} \]

That is what a matrix looks like. It is not what a matrix is, and the gap
between those two sentences is most of the trouble people have with the
subject.
\end{fr}

\begin{fr}
A matrix takes a vector and returns a vector. The one above doubles the first
component and triples the second, so it sends $(1, 1)$ to $(2, 3)$.

Two properties are what make it worth a name of its own. Apply it to a sum, and
apply it to a scaled vector, and say what you get in each case.
\dotline
\nextframe
\end{fr}

\begin{fr}
\begin{ansblock}
\[ A(u + v) = Au + Av \qquad A(cv) = c\,(Av) \]
\end{ansblock}

A function with those two properties is \emph{linear}. Check them on $A$ if you
have not: doubling and tripling componentwise cannot care whether you added the
vectors first, and it cannot care whether you scaled them first.

\textbf{Every single thing in this program is a consequence of those two
lines.} Not most of it \dash{} all of it. The rectangle of numbers is a way to
write such a function down, and it turns out that every linear function from
one finite-dimensional space to another can be written down that way.
\end{fr}

\mermaidfig{p06-matrix-is-a-function}{What a matrix looks like, and what it
is.}{a matrix is a function}

\begin{fr}
Which makes a test available that people usually apply by feel.

Is $f(x) = x + 3$ linear? It is a straight line, and the word \enquote{linear}
comes from \enquote{line}, so the answer is presumably yes.
\nextframe
\end{fr}

\begin{fr}
\ans{No.}
Put $x = 0$: $f(0) = 3$. But linearity gives $f(0) = f(0 \cdot 0) = 0 \cdot
f(0) = 0$, so a linear function must send zero to zero. One counterexample and
it is not linear.

\begin{trapbox}
\textbf{\enquote{Linear} in this book means the two properties, not
\enquote{a straight line}.} The English word and the mathematical word have
drifted apart, and $y = mx + c$ \dash{} the most quoted straight line there is
\dash{} is linear only when $c = 0$.

The general thing is called \emph{affine}: a linear map plus a shift. Every
layer in a network is affine rather than linear, because of the bias, and
almost every sentence in this program is about the linear part with the shift
carried along beside it.
\end{trapbox}
\end{fr}

\begin{fr}
The shift is worth one more frame, because it is why the collapse in section 5
needs a second line of algebra rather than one.

A layer computes $Wx + b$. The $W$ part is linear; the $+ b$ is not, and it is
the reason $f(0) \neq 0$ for a real layer. Program~\ref{prog:F06} built exactly
this object in one dimension, $y = wx + b$, and said the weights become a
matrix here.

So the objects line up: $w$ becomes $W$, $x$ becomes a vector, $b$ becomes a
vector, and the arithmetic that was multiplication becomes something that
needs its own definition.
\end{fr}

\begin{fr}
Before that definition, one thing worth noticing about what linearity buys,
because it is the reason the subject is tractable at all.

You know what $A$ does to $(1, 0)$ and to $(0, 1)$. From those two answers
alone, can you say what $A$ does to $(5, -2)$?
\nextframe
\end{fr}

\begin{fr}
\begin{ansblock}
Yes: $A(5, -2) = 5\,A(1, 0) - 2\,A(0, 1)$, by the two properties and nothing
else.
\end{ansblock}

\textbf{A linear map is completely determined by what it does to a basis.}
That is why a rectangle of numbers is enough to specify one: the columns of $A$
are exactly the images of the basis vectors, and everything else is a linear
combination Program~\ref{prog:P04} taught you to build.

An infinite amount of behaviour, pinned down by $n$ answers. Nothing else in
this book compresses that well.
\end{fr}

% ============================================================
\section{Multiplication is composition}
\label{sec:P06-composition}
\index{matrix!multiplication}

\begin{fr}
Now the definition everybody is taught as a procedure, and which is not one.

You have two linear maps: $B$ takes a vector and returns one, and $A$ takes
that and returns another. Doing both in turn is itself a function. Is it
linear?
\nextframe
\end{fr}

\begin{fr}
\begin{ansblock}
Yes, and the check is two lines: $A(B(u+v)) = A(Bu + Bv) = ABu + ABv$, and the
same for scaling.
\end{ansblock}

So the composition is a linear map, so it has a matrix of its own, and
\textbf{that matrix is what $AB$ means}. The definition is not
\enquote{multiply rows into columns} \dash{} that is the arithmetic that comes
out of it.

Read $AB$ as \emph{do $B$, then $A$}, right to left, the way $f(g(x))$ is read.
The order looks backwards because the notation puts the function on the left of
the thing it eats.
\end{fr}

\begin{fr}
Everything awkward about the procedure falls out of that reading, so it is
worth turning the first one over yourself.

The rule says the $i$-th row of $A$ meets the $j$-th column of $B$ and pairs
them off term by term. Nothing in \emph{do $B$, then $A$} obviously says that.
Where does it come from?
\dotline
\nextframe
\end{fr}

\begin{fr}
\begin{ansblock}
Entry $(i, j)$ is what the composed map does to the $j$-th basis vector, read
in the $i$-th coordinate.
\end{ansblock}

$B$ sends the $j$-th basis vector to its own $j$-th column, and $A$ then reads
that column a row at a time. The procedure is a consequence, and if you ever
forget it you can rebuild it from the sentence.

Which brings the property nobody has to memorise once it is a composition. Is
$(AB)C$ the same as $A(BC)$?
\nextframe
\end{fr}

\begin{fr}
\begin{ansblock}
Yes, necessarily. Both mean \emph{do $C$, then $B$, then $A$}.
\end{ansblock}

Associativity is not a fact about the arithmetic that somebody checked. It is
the observation that composing functions has never had a bracketing at all:
$A(B(C(x)))$ has one reading, and the brackets in $ABC$ are notation the
functions do not see.

The script confirms it to better than $\val{p06.assoc.err}$ on random
matrices, which is
worth doing once for the same reason the projection sweep was worth doing: an
identity agreeing with itself is not evidence, and here the arithmetic could in
principle have contradicted the theorem.
\end{fr}

\begin{fr}
\begin{note}
That is the whole reason this section is about composition rather than about
the procedure.

Associativity is not obvious from the rows-into-columns rule \dash{} written
out, it is a statement about a triple sum being reorderable, and it takes real
work to see. Read as composition it takes no work at all, because there was
never anything to prove.

The same will be true of everything else in this program. When a fact about
matrices looks arbitrary, the reading has slipped back to the rectangle.
\end{note}
\end{fr}

% ============================================================
\section{Why the order matters}
\label{sec:P06-order}
\index{matrix!non-commutativity}

\begin{fr}
Here is the property that one dimension cannot show you, and the reason is
worth saying before the example rather than after.

Program~\ref{prog:F05} and Program~\ref{prog:F06} both worked with $wx + b$,
where $w$ is a number. Two numbers multiplied in either order give the same
answer, always. So every intuition you have about composing those layers was
formed in a setting where \textbf{this section's whole subject is
invisible}.

Take two operations on the plane: a quarter turn anticlockwise, and flattening
everything onto the horizontal axis. Apply both, in both orders, to the point
$(1, 0)$.
\dotline
\nextframe
\end{fr}

\begin{fr}
\begin{ansblock}
Project then rotate: $(\val{p06.order.rp.x}, \val{p06.order.rp.y})$.
Rotate then project: $(\val{p06.order.pr.x}, \val{p06.order.pr.y})$.
\end{ansblock}

Follow it through. $(1, 0)$ already lies on the horizontal axis, so projecting
leaves it alone, and the quarter turn then lifts it to
$(\val{p06.order.rp.x}, \val{p06.order.rp.y})$. The other way, the turn sends
it to $(\val{p06.order.rp.x}, \val{p06.order.rp.y})$ first, and flattening kills
it entirely.

\textbf{One order loses the point and the other does not.} That is as far apart
as two answers get.
\end{fr}

\begin{fr}
\transcript{p06-order}

The two products are different matrices, and they are different because they
are different functions. There is no arithmetic accident here to be tidied up.

\begin{trapbox}
\textbf{$AB$ and $BA$ are not two ways of writing one thing, and the habit that
says otherwise came from arithmetic.}

Every multiplication you did before this book commuted. Numbers commute, and
because a one-dimensional layer's weight is a number, so does every product in
Programs~\ref{prog:F05} and~\ref{prog:F06}. The first time it stops being true
is the first time the weight becomes a matrix, which is here.

Every multiplication before this one commuted, so the habit is not a careless
one. It is a generalisation from the only case anybody had.
\end{trapbox}

The script draws $\val{p06.commute.trials}$ random pairs of $3 \times 3$
matrices and counts the ones that commute. Write down the number you expect
before reading on.
\dotline
\nextframe
\end{fr}

\begin{fr}
\ans{$\val{p06.commute.found}$.}
Not \enquote{few} \dash{} none, and it could not have been otherwise:
commuting is a condition on a set of measure zero, so a random draw misses it
with probability one. A pair that commutes has to be built, and every one you
will meet was.

\begin{aibox}
This is why \enquote{attention then the feed-forward block} and
\enquote{the feed-forward block then attention} describe two different models
rather than one model drawn two ways, and why a residual stream's order of
operations is part of the architecture rather than a diagram convention.

It is also why the transpose rule reverses: $(AB)\T = B\T A\T$. Reading it as
composition, the transpose runs the maps the other way, so the order has to
turn round with it \dash{} which is a sentence rather than an identity to
memorise.
\end{aibox}
\end{fr}

\mermaidfig{p06-order-matters}{Why two maps composed the other way round are a
different function.}{why order matters}

\begin{fr}
One clarification, because \enquote{does not commute} gets over-read.

It does not mean $AB \neq BA$ always. Some pairs do commute: any matrix with
the identity, any matrix with itself, two rotations of the plane about the same
centre.

Rotations in the plane commuting is worth holding on to, because rotations in
three dimensions do not \dash{} turn a book face-up-then-sideways and
sideways-then-face-up and the covers point different ways. \textbf{Two
dimensions is the special case, and it is the one everybody's intuition was
built in.}
\end{fr}

% ============================================================
\section{Shape is a type signature}
\label{sec:P06-shape}
\index{matrix!shape}

\begin{fr}
A linear map goes \emph{from} somewhere \emph{to} somewhere, and a matrix
carries both of those facts in its shape.

An $m \times n$ matrix takes a vector of $n$ components and returns one of $m$.
So the shape is not bookkeeping about how the numbers are stored \dash{} it is
the domain and the codomain, written down.

Given that, say which of $AB$ and $BA$ is even defined when $A$ is
$3 \times 5$ and $B$ is $5 \times 2$.
\nextframe
\end{fr}

\begin{fr}
\ans{$AB$ only, and it is $3 \times 2$.}
$B$ maps into a space of $5$ dimensions and $A$ reads one, so they meet. The
other order asks $B$ to read a space of $3$ dimensions when it expects $2$, and
there is nothing to compute.

\textbf{So a shape error is a type error.} The inner dimensions must match
because that is where one function's output meets the next one's input, and the
outer two survive because they are the ends of the composition.

\[ (m \times n)(n \times p) \longrightarrow (m \times p) \]
\end{fr}

\mermaidfig{p06-shape-is-a-type}{Why the inner dimensions have to agree, and
why the outer two survive.}{shape as a type signature}

\begin{fr}
\begin{aibox}
Which is why the error your framework raises names two numbers and looks
unhelpful. \code{mat1 and mat2 shapes cannot be multiplied (32x768 and
512x768)} is a type error reported by arithmetic, and the fix is never to
reshape until it stops complaining.

The question to ask is which of the two is wrong: has a layer been given the
wrong width, or has something been transposed that should not have been?
Those are different bugs with the same message, and the shapes tell you which
by saying what each factor claims to be a map \emph{between}.
\end{aibox}
\end{fr}

\begin{fr}
One shape that deserves its own frame, because it is where the reading pays
for itself.

A vector of $n$ components is an $n \times 1$ matrix. So $Ax$ is a matrix
product like any other, $(m \times n)(n \times 1) \to (m \times 1)$, and there
is no second rule for applying a matrix to a vector.

Which raises a question about a batch of them. You have $k$ vectors to push
through the same layer. What shape should they be stacked into?
\dotline
\nextframe
\end{fr}

\begin{fr}
\begin{ansblock}
$n \times k$ \dash{} one column each \dash{} and then $A$ times that is
$m \times k$, one output column each.
\end{ansblock}

The composition rule did all the work: nothing about $A$ changed and nothing
new was defined. Stacking $k$ inputs beside each other stacks $k$ outputs
beside each other, because each column is handled independently and linearity
never had anything to say about the other columns.

\begin{note}
Real frameworks stack the batch along the \emph{first} axis instead, so a
layer computes $XW\T$ with $X$ of shape $k \times n$. That is the same
composition with the convention flipped, and Program~\ref{prog:P07} is where
the axes get named properly \dash{} along with reshape, transpose and the
rank-4 arrays this program's two indices cannot describe.
\end{note}
\end{fr}

\begin{fr}
And the property that makes the shapes worth checking before running anything.

The composition $ABC$ is legal or it is not, and you can tell from three
shapes without touching a number. That is what a type signature is for: it
rejects a program before it runs, and the rejection names the join that does
not fit rather than the symptom fifty lines later.

\textbf{Read a stack of layers as a chain of shapes and most of the errors you
have hit become visible before the first batch.}
\end{fr}

% ============================================================
\section{The collapse, in more than one dimension}
\label{sec:P06-collapse}
\index{matrix!composition of layers}

\begin{fr}
Program~\ref{prog:F05} already made this argument and made it completely, in
one dimension: two hundred layers of $w x + b$ collapse to one, the extra
$199$ buy nothing but arithmetic, and depth on its own is not a capability.

Program~\ref{prog:F06} said the weights become a matrix here and asked whether
the argument survives. So take two layers, $W_1 x + b_1$ then
$W_2 h + b_2$, and compose them. Write down what comes out.
\dotline
\nextframe
\end{fr}

\begin{fr}
\begin{ansblock}
\[ W_2(W_1x + b_1) + b_2 = (W_2W_1)\,x + (W_2b_1 + b_2) \]
\end{ansblock}

One matrix and one bias vector. The argument survives exactly, and the two
steps that made it work are the two properties from section 1: $W_2$
distributes over the sum, and it carries a scalar straight through.

The script checks the identity over $\val{p06.collapse.trials}$ random inputs
at shapes $\val{p06.din} \to \val{p06.dmid} \to \val{p06.dout}$, worst
disagreement under $\val{p06.collapse.err}$ \dash{} which is rounding, not a
difference.
\end{fr}

\begin{fr}
Which is worth pausing on, because a stack of these is a large object.

Eight linear layers of width $\val{p06.stack.width}$ hold
$\val{p06.stack.params}$ parameters. The single layer they collapse to holds
$\val{p06.stack.collapsed}$.

So what fraction of that stack bought nothing?
\nextframe
\end{fr}

\begin{fr}
\ans{$\val{p06.stack.wasted.pct}\%$ of it \dash{} seven eighths.}
And the fraction is not a property of the width at all: $k$ layers collapse to
one, so $(k-1)/k$ of the parameters are redundant however wide they are. Add
layers and it approaches all of them.

\textbf{The waste is exactly the depth you thought you were buying.}
\end{fr}

\begin{fr}
Now the fix, which Program~\ref{prog:F05} named and this program can now say
something new about.

Put a $\relu$ between the two layers and the composition is
$W_2\,\relu(W_1x + b_1) + b_2$. F05's argument was that no amount of stretching
and sliding produces a corner, so the pair cannot be rewritten as a single
affine map.

In more than one dimension there is a cleaner way to see it, and it is worth
one frame because it needs no picture. An affine map sends any three collinear
points to three collinear points \dash{} it cannot bend a line. So: feed three
collinear inputs through the pair, and look at what comes out.
\nextframe
\end{fr}

\begin{fr}
\begin{ansblock}
They come out bent. Measured off the line by
$\val{p06.bend}$ of the spacing, against under $\val{p06.bend.affine}$ through
the collapsed matrix.
\end{ansblock}

The collapsed matrix keeps them collinear to rounding, because it must. The
$\relu$-separated pair does not, so \textbf{no single affine map reproduces it,
whatever that map might be} \dash{} the argument rules out all of them at once
rather than failing to find one.

That is F05's corner, restated in a form that generalises: linearity is a
constraint on what the function can do to a straight line, and the activation
is there to break it.
\end{fr}

\begin{fr}
\begin{aibox}
It also says what an activation is \emph{not} for, which is where the folklore
sits. Not squashing values into a range \dash{} $\relu$ is unbounded above and
does not. Not being biologically suggestive. Not smoothing anything.

It is there because the composition without it is provably a waste of
$\val{p06.stack.wasted.pct}\%$ of the parameters, and any function that breaks
linearity will do. That is why the field could swap the logistic for $\relu$
for GELU without rewriting the theory: the requirement is
\enquote{not affine}, and it is met with room to spare.
\end{aibox}
\end{fr}

\begin{fr}
The last thing composition gives you, and it costs nothing mathematically and a
great deal in practice.

$ABC$ is associative, so you may compute $(AB)C$ or $A(BC)$ and get the same
answer. Take $A$ as $\val{p06.dims.n} \times \val{p06.dims.m}$, $B$ as
$\val{p06.dims.m} \times \val{p06.dims.p}$ and $C$ as
$\val{p06.dims.p} \times \val{p06.dims.q}$, and count the multiplications each
way. An $(m \times n)(n \times p)$ product costs $mnp$ of them.
\dotline
\nextframe
\end{fr}

\begin{fr}
\begin{ansblock}
$(AB)C$ costs $\val{p06.cost.left}$ and $A(BC)$ costs
$\val{p06.cost.right}$ \dash{} a factor of $\val{p06.cost.ratio}$.
\end{ansblock}

Same answer, same theorem, and one bracketing does
$\val{p06.cost.ratio}$ times the work. $B$ has one column, so $AB$ is a thin
thing and cheap; $BC$ is the outer product of a column and a row, which is a
full $\val{p06.dims.n} \times \val{p06.dims.q}$ matrix, and then $A$ has to
multiply it.

\textbf{Associativity is free and bracketing is not.} Program~\ref{prog:P03}
built the vocabulary for exactly this: the two routes have identical results
and different bills.
\end{fr}

\begin{fr}
\begin{aibox}
That shape \dash{} a wide matrix, a very thin one, a wide one \dash{} is not a
contrived example. It is what a low-rank update looks like, and the reason such
an update is cheap is precisely the bracketing above: keep the thin factor in
the middle and never form the product of the outer two.

Program~\ref{prog:P08} is where the thin dimension gets its name and the
assumption behind it gets stated. What this program can say is the mechanical
half: the saving is associativity spent deliberately, and a framework that
formed $BC$ first would hand back the same numbers and
$\val{p06.cost.ratio}$ times the bill.
\end{aibox}
\end{fr}

\begin{fr}
The last frame, and it is about what the object bought.

Four sections ago a matrix was a rectangle of numbers with a multiplication
rule that had to be memorised. It is a function that respects addition and
scaling, and every one of the awkward facts came out of that sentence without
being separately learned: the product is composition, so it is associative and
does not commute; the shape is the two spaces, so a shape error is a type
error; and a stack of them is one of them, so the activation has a job.

Program~\ref{prog:P07} takes the two indices this program has and asks what
happens when there are four, which is where most people's mental model of an
array stops being a picture and has to become notation.
\end{fr}

% ============================================================
\begin{summarybox}
\sumitem{1--3}{\textbf{A matrix is a function, not a rectangle.} It takes a
  vector and returns one, and it is \emph{linear}: $A(u+v) = Au + Av$ and
  $A(cv) = c\,(Av)$. Everything else in the program is a consequence of those
  two lines, and every finite-dimensional linear map can be written as a
  matrix.}
\sumitem{4--6}{\textbf{\enquote{Linear} means those two properties, not
  \enquote{a straight line}.} $f(x) = x + 3$ is not linear, because a linear
  map must send zero to zero. A layer is \emph{affine} \dash{} linear plus the
  bias \dash{} and Program~\ref{prog:F06}'s $y = wx + b$ is the same object one
  dimension down.}
\sumitem{7--8}{\textbf{A linear map is determined by what it does to a basis},
  which is why a rectangle of numbers suffices to specify one: the columns are
  the images of the basis vectors and everything else is a linear combination.}
\sumitem{9--11}{\textbf{The product $AB$ is the composition} \dash{} do $B$,
  then $A$. The rows-into-columns procedure is a consequence of that, not the
  definition, and it can be rebuilt from the sentence whenever it is
  forgotten.}
\sumitem{12--14}{\textbf{Associativity is free}, because composing functions
  never had a bracketing: $A(B(C(x)))$ has one reading. Confirmed to better
  than $\val{p06.assoc.err}$, which is worth doing once because an identity
  agreeing with itself is not evidence.}
\sumitem{15--17}{\textbf{The order matters, and one dimension cannot show it},
  because numbers commute. A quarter turn and a flattening onto the horizontal
  axis send $(1, 0)$ to $(\val{p06.order.rp.x}, \val{p06.order.rp.y})$ one way
  and $(\val{p06.order.pr.x}, \val{p06.order.pr.y})$ the other \dash{} one
  order loses the point entirely.}
\sumitem{17--19}{Of $\val{p06.commute.trials}$ random $3 \times 3$ pairs,
  $\val{p06.commute.found}$ commute, and it could not be otherwise: commuting
  is a measure-zero condition. Some pairs do \dash{} rotations of the plane
  about one centre \dash{} and \textbf{rotations in three dimensions do not},
  which is why two dimensions is the misleading case.}
\sumitem{20--22}{\textbf{Shape is the domain and the codomain written down}, so
  $(m \times n)(n \times p) \to (m \times p)$: the inner dimensions meet where
  one function's output is the next one's input. A shape error is a type error,
  and the shapes say which of the two plausible bugs it is.}
\sumitem{23--25}{A vector is an $n \times 1$ matrix, so $Ax$ needs no second
  rule, and a batch of $k$ of them is $n \times k$ \dash{} \textbf{linearity
  never had anything to say about the other columns}. The axis conventions and
  the rank-4 case are Program~\ref{prog:P07}'s.}
\sumitem{26--29}{\textbf{Two affine layers are one}: $W_2(W_1x + b_1) + b_2 =
  (W_2W_1)x + (W_2b_1 + b_2)$, checked over $\val{p06.collapse.trials}$ inputs
  to better than $\val{p06.collapse.err}$. $k$ layers waste $(k-1)/k$ of their parameters
  whatever their width \dash{} $\val{p06.stack.wasted.pct}\%$ at eight layers
  \dash{} which is exactly the depth you thought you were buying.}
\sumitem{30--32}{\textbf{An affine map cannot bend a straight line}, so three
  collinear inputs come out collinear (under $\val{p06.bend.affine}$) through
  the collapsed matrix and bent ($\val{p06.bend}$) through the
  $\relu$-separated pair. That rules out every single affine map at once, and it is why the
  requirement on an activation is only \enquote{not affine}.}
\sumitem{33--34}{\textbf{Associativity is free and bracketing is not}: the same
  triple product costs $\val{p06.cost.left}$ multiplications one way and
  $\val{p06.cost.right}$ the other, a factor of $\val{p06.cost.ratio}$. Keeping
  a thin factor in the middle is that saving spent deliberately.}
\end{summarybox}

% ============================================================
\begin{testexercises}
\item State the two properties that make a function linear.
  \answerto{$f(u + v) = f(u) + f(v)$ and $f(cv) = c\,f(v)$. Together they force
  $f(0) = 0$, which is the quickest test there is.}
\item Say whether $f(x, y) = (2x, x + y, 0)$ is linear, with the reason.
  \answerto{Yes. Every output component is a sum of input components times
  constants, with no constant term, so both properties hold. Its matrix is
  $3 \times 2$.}
\item $A$ is $4 \times 7$ and $B$ is $7 \times 2$. Give the shape of $AB$, and
  say whether $BA$ exists.
  \answerto{$AB$ is $4 \times 2$. $BA$ does not exist: $B$ maps into a space of
  $7$ dimensions and $A$ reads $7$, but the other order asks $B$ to read $4$
  when it expects $2$.}
\item Say what $AB$ means when $A$ and $B$ are read as functions, and use it to
  explain why $(AB)C = A(BC)$ without computing anything.
  \answerto{$AB$ is \emph{do $B$, then $A$}. Both bracketings mean \emph{do
  $C$, then $B$, then $A$}, and composing functions has no bracketing to get
  wrong, so there is nothing to prove.}
\item Give two maps of the plane whose two orders disagree, and say where a
  chosen point goes each way.
  \answerto{A quarter turn and a projection onto the horizontal axis. On
  $(1, 0)$: project then rotate gives $(\val{p06.order.rp.x},
  \val{p06.order.rp.y})$, rotate then project gives $(\val{p06.order.pr.x},
  \val{p06.order.pr.y})$.}
\item Say why the failure of $AB = BA$ could not have been noticed in
  Programs~\ref{prog:F05} and~\ref{prog:F06}.
  \answerto{Their weights are numbers, and numbers commute. The phenomenon
  appears the moment the weight becomes a matrix, which is the first place
  composition stops being multiplication.}
\item Collapse $W_2(W_1x + b_1) + b_2$ into a single affine map.
  \answerto{$(W_2W_1)x + (W_2b_1 + b_2)$. One matrix and one bias vector, by
  distributing $W_2$ over the sum and passing it through the scalar.}
\item A stack of $k$ linear layers of width $d$ has no activations. Say what
  fraction of its parameters is redundant, and whether the width matters.
  \answerto{$(k-1)/k$, and the width does not matter at all \dash{} at
  $k = \val{p06.stack.layers}$ it is $\val{p06.stack.wasted.pct}\%$ whatever
  $d$ is. The waste is exactly the depth.}
\item Say why three collinear inputs settle whether a pair of layers is affine.
  \answerto{An affine map sends collinear points to collinear points, always.
  If the outputs are not collinear then no affine map produces them, which
  rules out every candidate at once rather than failing to find one.}
\item Two bracketings of one triple product cost $\val{p06.cost.left}$ and
  $\val{p06.cost.right}$ multiplications. Say what that does and does not tell
  you.
  \answerto{It says nothing about the answer, which is identical because
  associativity is a theorem. It says the bill differs by a factor of
  $\val{p06.cost.ratio}$, so bracketing is an engineering decision that the
  mathematics leaves entirely free.}
\end{testexercises}

% ============================================================
\begin{furtherproblems}
\item Show that the columns of $A$ are the images of the basis vectors.
  \answerto{$e_j$ has a $1$ in position $j$ and zeros elsewhere, so $Ae_j$
  picks out column $j$ of $A$ by the definition of the product. That is what
  makes a rectangle of numbers enough to specify a map.}
\item Show that a composition of linear maps is linear, from the definition.
  \answerto{$A(B(u+v)) = A(Bu + Bv) = ABu + ABv$, and $A(B(cv)) = A(c\,Bv) =
  c\,ABv$. Both steps use each map's own linearity once.}
\item Give a pair of matrices that do commute and say what is special about
  them.
  \answerto{Any matrix and the identity; any matrix and itself; two rotations
  of the plane about the same centre. In the last case both are rotations by an
  angle and the two orders add the same two angles.}
\item Show that $(AB)\T = B\T A\T$ follows from reading the product as a
  composition.
  \answerto{The transpose runs the maps the other way round, so the composition
  has to reverse with it. Written out, entry $(i,j)$ of the left side is row
  $j$ of $A$ against column $i$ of $B$, which is exactly entry $(i,j)$ of the
  right.}
\item A layer is $W x + b$ with $W$ of shape $m \times n$. Say what shape a
  batch of $k$ inputs has under this program's convention, and what shape comes
  out.
  \answerto{$n \times k$ in and $m \times k$ out, one column each. Nothing new
  is defined \dash{} the composition rule already handles it, because each
  column is independent of the others.}
\item Explain why the fraction of wasted parameters in a linear stack does not
  depend on the width.
  \answerto{$k$ layers of width $d$ hold $k(d^{2} + d)$ parameters and collapse
  to $d^{2} + d$. The ratio is $1/k$ and $d$ cancels, so the waste is
  $(k-1)/k$ at every width.}
\item Say why \enquote{the activation squashes values into a range} cannot be
  the reason activations exist.
  \answerto{$\relu$ is unbounded above and does not squash anything into a
  range, and it works. The requirement met by every activation in use is only
  that it is not affine, because an affine one would let the layers collapse.}
\item Given $A$ of shape $d \times d$, $B$ of shape $d \times r$ and $C$ of
  shape $r \times d$ with $r$ small, say which bracketing of $ABC$ to prefer
  and by roughly how much.
  \answerto{$(AB)C$: it costs $d^{2}r + d r d = 2d^{2}r$ against $A(BC)$'s
  $d r d + d^{3} = d^{2}r + d^{3}$. For small $r$ the second is dominated by
  $d^{3}$, so the saving is about $d/(2r)$ \dash{} which is why a low-rank
  update is cheap only if you never form the outer product.}
\end{furtherproblems}
